Точку $z \in \Comp$ можно построить тогда и только тогда, когда можно построить точки $Re\ z, Im\ z$.

1) Если задана точка $z \in \Comp$, то можно построить проекции этой
точки на обе оси (задача сводится к построению перпендикуляра к прямой
через заданную точку). Эти проекции и будут $Re\ z, Im\ z$.

2) Если нам известны $Re\ z$ и $Im\ z$ (они отмечены на осях), то
можно восстановить перпендикуляры из этих точек к соответствующим осям.
Точка пересечения этих перпендикуляров и будет искомой точкой $z$.